\begin{itemize}
    \item Proper rotation by an angle $\varphi$ with the center of rotation at the origin: \\ $x' = x \cos \varphi - y \sin \varphi \\ 
    y' = x \sin \varphi + y \cos \varphi$
    \item Distance between points on a sphere: 
    $L = R \cdot \arccos(\cos \theta_{1} \cdot \cos \theta_{2} + \sin \theta_{1} \cdot \sin \theta_{2} \cdot \cos(\varphi_{1} - \varphi_{2}))$
    where $\theta$ – latitudes (from $-\pi$ to $\pi$), $\varphi$ – longitudes (from $-\pi$ to $\pi$)
    \item Volume of a spherical cap: $V = \pi h^{2}(R - \frac{1}{3}h)$, where $h$ -- height from the apex of the sector to the cutting plane
    \item Surface area of a spherical cap: $S = 2\pi Rh$, where $h$ -- height
    % \item $2^{23} \cdot 7 \cdot 17 + 1 = 998.244.353$ --- prime, primitive root -- $3$
    \item Gray code: $g_{n} = n \oplus \frac{n}{2}$
    \item Fibonacci numbers: \\ $\displaystyle F_{0} = 0, F_{1} = 1, F_{n} = \frac{(\frac{1 + \sqrt{5}}{2})^{n} - (\frac{1 - \sqrt{5}}{2})^{n}}{\sqrt{5}}$
    \item Sum-xor property: $a + b = a \oplus b + 2(a\&b), a + b = a|b + a\&b, a \oplus b = a|b - a\&b$
    \item Number of faces in a planar graph(including the infinite one): $R = 2 - V + E$
    \item Sum of an arithmetic progression: $S_{n} = \frac{n(a_{1} + a_{n})}{2}$
    \item Sum of a geometric progression: $S_{n} = \frac{b_{1}(q^{n} - 1)}{q - 1}$
    \item Determinants of matrices
    $$
    \left|
    \begin{array}{c c}
        a & b \\
        c & d
    \end{array}
    \right| = ad - bc
    $$
    \begin{multline*}
    \left|
    \begin{array}{c c c}
        a_{1} & b_{1} & c_{1} \\
        a_{2} & b_{2} & c_{2} \\
        a_{3} & b_{3} & c_{3} \\
    \end{array}
    \right| = a_{1}b_{2}c_{3} + a_{3}b_{1}c_{2} + a_{2}b_{3}c_{1} -\\- a_{3}b_{2}c_{1} - a_{1}b_{3}c_{2} - a_{2}b_{1}c_{3}
    \end{multline*}
    $\Delta = \sum_{j = 1}^{n} (-1)^{j + 1} \cdot a_{1, j} \cdot \bar{M_{j}^{1}}$, $\bar{M_{j}^{1}}$ --- determinant of the matrix obtained by deleting row $1$ and column $j$.
    \item \textbf{Cramer's rule.} $\det A \neq 0 \implies$ unique solution. Otherwise $0$ or $\infty$.
    Solutions: $\displaystyle x_{i} = \frac{\Delta_{i}}{\Delta}$. In $\Delta_{i}$, the column of coefficients of the corresponding unknown is replaced by the column of constant terms of the system.
\end{itemize}
